\medskip (i) Orient the distinguished curves. The global
orientation will be denoted $\twoheadrightarrow\!\!-$ and the tangle
orientation by $\rightarrow\!\!-$.

\smallskip (ii) At a point of intersection between the
distinguished curves and the strings of the tangle, insert 2-boxes
containing $u$ or $u^{-1}$ according to the following conventions
\[
	\begin{picture}(0,0)%
\epsfig{file=ps/pic45.ps}%
\end{picture}%
\setlength{\unitlength}{0.00041700in}%
%
\begingroup\makeatletter\ifx\SetFigFont\undefined
% extract first six characters in \fmtname
\def\x#1#2#3#4#5#6#7\relax{\def\x{#1#2#3#4#5#6}}%
\expandafter\x\fmtname xxxxxx\relax \def\y{splain}%
\ifx\x\y   % LaTeX or SliTeX?
\gdef\SetFigFont#1#2#3{%
  \ifnum #1<17\tiny\else \ifnum #1<20\small\else
  \ifnum #1<24\normalsize\else \ifnum #1<29\large\else
  \ifnum #1<34\Large\else \ifnum #1<41\LARGE\else
     \huge\fi\fi\fi\fi\fi\fi
  \csname #3\endcsname}%
\else
\gdef\SetFigFont#1#2#3{\begingroup
  \count@#1\relax \ifnum 25<\count@\count@25\fi
  \def\x{\endgroup\@setsize\SetFigFont{#2pt}}%
  \expandafter\x
    \csname \romannumeral\the\count@ pt\expandafter\endcsname
    \csname @\romannumeral\the\count@ pt\endcsname
  \csname #3\endcsname}%
\fi
\fi\endgroup
\begin{picture}(5497,4270)(579,-4684)
\put(4651,-1486){\makebox(0,0)[lb]{\smash{\SetFigFont{12}{14.4}{rm}$u$}}}
\put(5026,-3436){\makebox(0,0)[lb]{\smash{\special{ps:gsave currentpoint currentpoint translate
-180.0 rotate neg exch neg exch translate}\SetFigFont{12}{14.4}{rm}$u^{-1}$\special{ps:currentpoint grestore moveto}}}}
\put(6076,-1036){\makebox(0,0)[lb]{\smash{\SetFigFont{12}{14.4}{rm}(global orientation)}}}
\put(6076,-2215){\makebox(0,0)[lb]{\smash{\SetFigFont{12}{14.4}{rm}going out)}}}
\put(6076,-3436){\makebox(0,0)[lb]{\smash{\SetFigFont{12}{14.4}{rm}(global orientation)}}}
\put(6076,-1711){\makebox(0,0)[lb]{\smash{\SetFigFont{12}{14.4}{rm}(=tangle orientation}}}
\put(6076,-4111){\makebox(0,0)[lb]{\smash{\SetFigFont{12}{14.4}{rm}(=tangel orientation}}}
\put(6076,-4615){\makebox(0,0)[lb]{\smash{\SetFigFont{12}{14.4}{rm}going in)}}}
\put(4876,-4036){\makebox(0,0)[lb]{\smash{\SetFigFont{12}{14.4}{rm}$*$}}}
\put(4351,-1261){\makebox(0,0)[lb]{\smash{\SetFigFont{12}{14.4}{rm}$*$}}}
\end{picture}

\]
\bs
\ni Thus along a distinguished curve one alternately meets $u$ or
$u^{-1}$.

From now on we suppose for convenience that $\D = 1$.

\bs\ni
{\bf Lemma 2.11.2} The Reidemeister type  II moves are
satisfied, i.e.,
\[
	\begin{picture}(0,0)%
\epsfig{file=ps/pic46.ps}%
\end{picture}%
\setlength{\unitlength}{0.00041700in}%
%
\begingroup\makeatletter\ifx\SetFigFont\undefined
% extract first six characters in \fmtname
\def\x#1#2#3#4#5#6#7\relax{\def\x{#1#2#3#4#5#6}}%
\expandafter\x\fmtname xxxxxx\relax \def\y{splain}%
\ifx\x\y   % LaTeX or SliTeX?
\gdef\SetFigFont#1#2#3{%
  \ifnum #1<17\tiny\else \ifnum #1<20\small\else
  \ifnum #1<24\normalsize\else \ifnum #1<29\large\else
  \ifnum #1<34\Large\else \ifnum #1<41\LARGE\else
     \huge\fi\fi\fi\fi\fi\fi
  \csname #3\endcsname}%
\else
\gdef\SetFigFont#1#2#3{\begingroup
  \count@#1\relax \ifnum 25<\count@\count@25\fi
  \def\x{\endgroup\@setsize\SetFigFont{#2pt}}%
  \expandafter\x
    \csname \romannumeral\the\count@ pt\expandafter\endcsname
    \csname @\romannumeral\the\count@ pt\endcsname
  \csname #3\endcsname}%
\fi
\fi\endgroup
\begin{picture}(8744,1544)(879,-2183)
\put(5026,-1411){\makebox(0,0)[lb]{\smash{\SetFigFont{12}{14.4}{rm},}}}
\put(2776,-1486){\makebox(0,0)[lb]{\smash{\SetFigFont{12}{14.4}{rm}$=$}}}
\put(7561,-1456){\makebox(0,0)[lb]{\smash{\SetFigFont{12}{14.4}{rm}$=$}}}
\end{picture}

\]
\ni where either of the two curves is distinguished and its global
orientation is arbitrary.

\bs
{\sc Proof.} This is just a re-expression of bi-invertibility. \qed

\bs\ni
{\bf Theorem 2.11.3} If $P$ is a general planar algebra and
$u$ is bi-invertible, let $P^u_k$ be
\bs
\[
\{x\in P_k\ \text{s.t.}\begin{picture}(0,0)%
\epsfig{file=ps/pic47.ps}%
\end{picture}%
\setlength{\unitlength}{0.00041700in}%
%
\begingroup\makeatletter\ifx\SetFigFont\undefined
% extract first six characters in \fmtname
\def\x#1#2#3#4#5#6#7\relax{\def\x{#1#2#3#4#5#6}}%
\expandafter\x\fmtname xxxxxx\relax \def\y{splain}%
\ifx\x\y   % LaTeX or SliTeX?
\gdef\SetFigFont#1#2#3{%
  \ifnum #1<17\tiny\else \ifnum #1<20\small\else
  \ifnum #1<24\normalsize\else \ifnum #1<29\large\else
  \ifnum #1<34\Large\else \ifnum #1<41\LARGE\else
     \huge\fi\fi\fi\fi\fi\fi
  \csname #3\endcsname}%
\else
\gdef\SetFigFont#1#2#3{\begingroup
  \count@#1\relax \ifnum 25<\count@\count@25\fi
  \def\x{\endgroup\@setsize\SetFigFont{#2pt}}%
  \expandafter\x
    \csname \romannumeral\the\count@ pt\expandafter\endcsname
    \csname @\romannumeral\the\count@ pt\endcsname
  \csname #3\endcsname}%
\fi
\fi\endgroup
\begin{picture}(6929,3059)(879,-3083)
\put(4216,-1591){\makebox(0,0)[lb]{\smash{\SetFigFont{12}{14.4}{rm}$=$}}}
\put(2356,-1876){\makebox(0,0)[lb]{\smash{\SetFigFont{12}{14.4}{rm}$x$}}}
\put(6571,-1246){\makebox(0,0)[lb]{\smash{\SetFigFont{12}{14.4}{rm}$y$}}}
\end{picture}
 
\mbox{\ \ for some}\ \ y\in  P_k\}.
\]
\bs
Then $P^u$ is a  general planar subalgebra, $P^u$ is a (general)
planar $*$-subalgebra if $P$ is a planar $*$-algebra and $u$ is
bi-unitary. The properties of being planar, $C^*$ and spherical
are inherited from $P$.

\bs
{\sc Proof.} That $P^u$ is a subalgebra is obvious. The
$*$-property is more interesting. Applying $*$ to the pictures we
obtain (using $u^*=u^{-1}$)
\[
	\begin{picture}(0,0)%
\epsfig{file=ps/pic48.ps}%
\end{picture}%
\setlength{\unitlength}{0.00041700in}%
%
\begingroup\makeatletter\ifx\SetFigFont\undefined
% extract first six characters in \fmtname
\def\x#1#2#3#4#5#6#7\relax{\def\x{#1#2#3#4#5#6}}%
\expandafter\x\fmtname xxxxxx\relax \def\y{splain}%
\ifx\x\y   % LaTeX or SliTeX?
\gdef\SetFigFont#1#2#3{%
  \ifnum #1<17\tiny\else \ifnum #1<20\small\else
  \ifnum #1<24\normalsize\else \ifnum #1<29\large\else
  \ifnum #1<34\Large\else \ifnum #1<41\LARGE\else
     \huge\fi\fi\fi\fi\fi\fi
  \csname #3\endcsname}%
\else
\gdef\SetFigFont#1#2#3{\begingroup
  \count@#1\relax \ifnum 25<\count@\count@25\fi
  \def\x{\endgroup\@setsize\SetFigFont{#2pt}}%
  \expandafter\x
    \csname \romannumeral\the\count@ pt\expandafter\endcsname
    \csname @\romannumeral\the\count@ pt\endcsname
  \csname #3\endcsname}%
\fi
\fi\endgroup
\begin{picture}(6974,3055)(865,-3084)
\put(4216,-1591){\makebox(0,0)[lb]{\smash{\SetFigFont{12}{14.4}{rm}$=$}}}
\put(2251,-1336){\makebox(0,0)[lb]{\smash{\SetFigFont{12}{14.4}{rm}$x^*$}}}
\put(6556,-1891){\makebox(0,0)[lb]{\smash{\SetFigFont{12}{14.4}{rm}$y^*$}}}
\end{picture}

\]
\ni Now we surround these pictures with the annular tangle
\[
	\begin{picture}(0,0)%
\epsfig{file=ps/pic49.ps}%
\end{picture}%
\setlength{\unitlength}{0.00041700in}%
%
\begingroup\makeatletter\ifx\SetFigFont\undefined
% extract first six characters in \fmtname
\def\x#1#2#3#4#5#6#7\relax{\def\x{#1#2#3#4#5#6}}%
\expandafter\x\fmtname xxxxxx\relax \def\y{splain}%
\ifx\x\y   % LaTeX or SliTeX?
\gdef\SetFigFont#1#2#3{%
  \ifnum #1<17\tiny\else \ifnum #1<20\small\else
  \ifnum #1<24\normalsize\else \ifnum #1<29\large\else
  \ifnum #1<34\Large\else \ifnum #1<41\LARGE\else
     \huge\fi\fi\fi\fi\fi\fi
  \csname #3\endcsname}%
\else
\gdef\SetFigFont#1#2#3{\begingroup
  \count@#1\relax \ifnum 25<\count@\count@25\fi
  \def\x{\endgroup\@setsize\SetFigFont{#2pt}}%
  \expandafter\x
    \csname \romannumeral\the\count@ pt\expandafter\endcsname
    \csname @\romannumeral\the\count@ pt\endcsname
  \csname #3\endcsname}%
\fi
\fi\endgroup
\begin{picture}(3044,3045)(3867,-3085)
\end{picture}

\]
\ni and then apply type II Reidemeister moves. We see that $x^*\in
P^u$ if $x$ does (though note that the ``$y$'' for $x^*$ is
$y^*$, but rotated.

To show that $P^u$ is a general planar subalgebra, we only have to
show by 1.18 that it is invariant under $\cal A(\phi)$. But if we
arrange the annular tangle $A$ so that all critical points of the
height function on strings are local maxima and minima, the
distinguished line
\[
	\begin{picture}(0,0)%
\epsfig{file=ps/pic50.ps}%
\end{picture}%
\setlength{\unitlength}{0.00041700in}%
%
\begingroup\makeatletter\ifx\SetFigFont\undefined
% extract first six characters in \fmtname
\def\x#1#2#3#4#5#6#7\relax{\def\x{#1#2#3#4#5#6}}%
\expandafter\x\fmtname xxxxxx\relax \def\y{splain}%
\ifx\x\y   % LaTeX or SliTeX?
\gdef\SetFigFont#1#2#3{%
  \ifnum #1<17\tiny\else \ifnum #1<20\small\else
  \ifnum #1<24\normalsize\else \ifnum #1<29\large\else
  \ifnum #1<34\Large\else \ifnum #1<41\LARGE\else
     \huge\fi\fi\fi\fi\fi\fi
  \csname #3\endcsname}%
\else
\gdef\SetFigFont#1#2#3{\begingroup
  \count@#1\relax \ifnum 25<\count@\count@25\fi
  \def\x{\endgroup\@setsize\SetFigFont{#2pt}}%
  \expandafter\x
    \csname \romannumeral\the\count@ pt\expandafter\endcsname
    \csname @\romannumeral\the\count@ pt\endcsname
  \csname #3\endcsname}%
\fi
\fi\endgroup
\begin{picture}(3644,3043)(3579,-3083)
\put(5026,-1636){\makebox(0,0)[lb]{\smash{\SetFigFont{12}{14.4}{rm}$\pi_A(x)$}}}
\end{picture}

\]
\ni can be moved through $\pi_A(x)$, close to \ $\tangle{x}$ 
using only planar isotopy and type II Reidemeister moves.
It can then go past $x$, producing a $y$, and down to the bottom
via type II Reidemeister moves. Thus  $\pi_A(x)$ is in $P^u$.

Planarity, positivity and sphericity are all
inherited. \qed

\bs {\bf Notes.} (i) We will see that $P^u$ may be planar even
when $P$ is only general planar.

(ii) When $k=2$ the equation of Theorem 2.11.3 is an abstract
version of the Yang-Baxter equation ([Ba]).

\bs We now recast the equation of Theorem 2.11.3 in some equivalent
forms which reveal some of its structure. The idea of the equation,
as seen clearly in the proof of 2.11.3 is just that the
distinguished lines can move freely past the boxes. As stated this
requires a special configuration as the distinguished line
approaches a box, but by Reidemeister type II invariance any
approach will do. We record this below, keeping the notation of
2.11.3. (Note that it is somewhat cumbersome to force all the
pictures to fit appropriately into the standard $k$-box. We use a
disk with $2k$ boundary points-as in the introduction.)

\medskip
{\bf Proposition 2.11.4} An $x$ in $P_k$ is in $P^u_k$ iff
there is a $y$ in $P^u_k$ with
\[
	\begin{picture}(0,0)%
\epsfig{file=ps/pic51.ps}%
\end{picture}%
\setlength{\unitlength}{0.00041700in}%
%
\begingroup\makeatletter\ifx\SetFigFont\undefined
% extract first six characters in \fmtname
\def\x#1#2#3#4#5#6#7\relax{\def\x{#1#2#3#4#5#6}}%
\expandafter\x\fmtname xxxxxx\relax \def\y{splain}%
\ifx\x\y   % LaTeX or SliTeX?
\gdef\SetFigFont#1#2#3{%
  \ifnum #1<17\tiny\else \ifnum #1<20\small\else
  \ifnum #1<24\normalsize\else \ifnum #1<29\large\else
  \ifnum #1<34\Large\else \ifnum #1<41\LARGE\else
     \huge\fi\fi\fi\fi\fi\fi
  \csname #3\endcsname}%
\else
\gdef\SetFigFont#1#2#3{\begingroup
  \count@#1\relax \ifnum 25<\count@\count@25\fi
  \def\x{\endgroup\@setsize\SetFigFont{#2pt}}%
  \expandafter\x
    \csname \romannumeral\the\count@ pt\expandafter\endcsname
    \csname @\romannumeral\the\count@ pt\endcsname
  \csname #3\endcsname}%
\fi
\fi\endgroup
\begin{picture}(7380,3750)(136,-3958)
\put(4111,-1921){\makebox(0,0)[lb]{\smash{\SetFigFont{12}{14.4}{rm}$=$}}}
\put(2521,-3886){\makebox(0,0)[lb]{\smash{\SetFigFont{12}{14.4}{rm}$q$ points}}}
\put(6766,-3856){\makebox(0,0)[lb]{\smash{\SetFigFont{12}{14.4}{rm}$q$ points}}}
\put(4381,-511){\makebox(0,0)[lb]{\smash{\SetFigFont{12}{14.4}{rm}$p$ points}}}
\put(136,-496){\makebox(0,0)[lb]{\smash{\SetFigFont{12}{14.4}{rm}$p$ points}}}
\put(2041,-1906){\makebox(0,0)[lb]{\smash{\SetFigFont{12}{14.4}{rm}$x$}}}
\put(6256,-1936){\makebox(0,0)[lb]{\smash{\SetFigFont{12}{14.4}{rm}$y$}}}
\end{picture}

\]
\ni where $p+q=2k$, where the pictures may be isotoped in any way
so that the annular region becomes the standard annular region
(with $*$ anywhere allowed by the orientations).

\medskip
{\sc Proof.} If $p>q$, surround the two pictures with the annular
tangle
\[
	\begin{picture}(0,0)%
\epsfig{file=ps/pic52.ps}%
\end{picture}%
\setlength{\unitlength}{0.00041700in}%
%
\begingroup\makeatletter\ifx\SetFigFont\undefined
% extract first six characters in \fmtname
\def\x#1#2#3#4#5#6#7\relax{\def\x{#1#2#3#4#5#6}}%
\expandafter\x\fmtname xxxxxx\relax \def\y{splain}%
\ifx\x\y   % LaTeX or SliTeX?
\gdef\SetFigFont#1#2#3{%
  \ifnum #1<17\tiny\else \ifnum #1<20\small\else
  \ifnum #1<24\normalsize\else \ifnum #1<29\large\else
  \ifnum #1<34\Large\else \ifnum #1<41\LARGE\else
     \huge\fi\fi\fi\fi\fi\fi
  \csname #3\endcsname}%
\else
\gdef\SetFigFont#1#2#3{\begingroup
  \count@#1\relax \ifnum 25<\count@\count@25\fi
  \def\x{\endgroup\@setsize\SetFigFont{#2pt}}%
  \expandafter\x
    \csname \romannumeral\the\count@ pt\expandafter\endcsname
    \csname @\romannumeral\the\count@ pt\endcsname
  \csname #3\endcsname}%
\fi
\fi\endgroup
\begin{picture}(3180,2890)(136,-3098)
\put(136,-496){\makebox(0,0)[lb]{\smash{\SetFigFont{12}{14.4}{rm}$k$ points}}}
\end{picture}

\]
\ni and  use isotopy, rotation (if necessary to get $*$ in the
right place), and type II Reidemeister moves to obtain the same
picture as in 2.11.3. \qed

\bs{\bf Definition 2.11.5.} Given $P$ and $u$ as above, define
$\s_u:P_k\to P_{k+1}$ by the tangle below:
\[
	\begin{picture}(0,0)%
\epsfig{file=ps/pic53.ps}%
\end{picture}%
\setlength{\unitlength}{0.00041700in}%
%
\begingroup\makeatletter\ifx\SetFigFont\undefined
% extract first six characters in \fmtname
\def\x#1#2#3#4#5#6#7\relax{\def\x{#1#2#3#4#5#6}}%
\expandafter\x\fmtname xxxxxx\relax \def\y{splain}%
\ifx\x\y   % LaTeX or SliTeX?
\gdef\SetFigFont#1#2#3{%
  \ifnum #1<17\tiny\else \ifnum #1<20\small\else
  \ifnum #1<24\normalsize\else \ifnum #1<29\large\else
  \ifnum #1<34\Large\else \ifnum #1<41\LARGE\else
     \huge\fi\fi\fi\fi\fi\fi
  \csname #3\endcsname}%
\else
\gdef\SetFigFont#1#2#3{\begingroup
  \count@#1\relax \ifnum 25<\count@\count@25\fi
  \def\x{\endgroup\@setsize\SetFigFont{#2pt}}%
  \expandafter\x
    \csname \romannumeral\the\count@ pt\expandafter\endcsname
    \csname @\romannumeral\the\count@ pt\endcsname
  \csname #3\endcsname}%
\fi
\fi\endgroup
\begin{picture}(4448,3053)(1576,-3082)
\put(4366,-1334){\makebox(0,0)[lb]{\smash{\SetFigFont{12}{14.4}{rm}$x$}}}
\put(1576,-1561){\makebox(0,0)[lb]{\smash{\SetFigFont{12}{14.4}{rm}$\sigma_u(x)=$}}}
\end{picture}

\]

\bs\ni
{\bf Proposition 2.11.6} The map $\s_u$ is a unital
endomorphism of the filtered algebra $P$ (a $*$-endomorphism if $u$
is unitary) and $P^u_k=\{x\mid \s_u(x)\in P_{1,k+1}\}$.

\bs
{\sc Proof.}  That $\s_u$ preserves multiplication follows from
type II Reidemeister moves. The alternative definition of $P^u$ is
just the case $q=0$, with $*$ appropriately placed, in 2.11.4. \qed

\bs Note that the endomorphism $\s_u$ is the obvious ``shift de un"
when restricted to the Temperley-Lieb subalgebra.

The condition of $2.11.5$ involved a pair $(x,y)$. In fact $x$ is
determined by $y$ and vice versa as we now record.

\bs\ni
{\bf Proposition 2.11.7} Suppose $P$ is a general planar
algebra. If $x\in P^u$ with \ $\s_u(x)=$ $\boxi{y}$.
Then

\medskip
\[
	\begin{picture}(0,0)%
\epsfig{file=ps/pic54.ps}%
\end{picture}%
\setlength{\unitlength}{0.00041700in}%
%
\begingroup\makeatletter\ifx\SetFigFont\undefined
% extract first six characters in \fmtname
\def\x#1#2#3#4#5#6#7\relax{\def\x{#1#2#3#4#5#6}}%
\expandafter\x\fmtname xxxxxx\relax \def\y{splain}%
\ifx\x\y   % LaTeX or SliTeX?
\gdef\SetFigFont#1#2#3{%
  \ifnum #1<17\tiny\else \ifnum #1<20\small\else
  \ifnum #1<24\normalsize\else \ifnum #1<29\large\else
  \ifnum #1<34\Large\else \ifnum #1<41\LARGE\else
     \huge\fi\fi\fi\fi\fi\fi
  \csname #3\endcsname}%
\else
\gdef\SetFigFont#1#2#3{\begingroup
  \count@#1\relax \ifnum 25<\count@\count@25\fi
  \def\x{\endgroup\@setsize\SetFigFont{#2pt}}%
  \expandafter\x
    \csname \romannumeral\the\count@ pt\expandafter\endcsname
    \csname @\romannumeral\the\count@ pt\endcsname
  \csname #3\endcsname}%
\fi
\fi\endgroup
\begin{picture}(10148,3055)(451,-3088)
\put(5401,-1711){\makebox(0,0)[lb]{\smash{\SetFigFont{12}{14.4}{rm}and }}}
\put(451,-1711){\makebox(0,0)[lb]{\smash{\SetFigFont{12}{14.4}{rm}$y$}}}
\put(976,-1711){\makebox(0,0)[lb]{\smash{\SetFigFont{12}{14.4}{rm}$=$}}}
\put(1726,-1711){\makebox(0,0)[lb]{\smash{\SetFigFont{12}{14.4}{rm}$\frac{1}{\delta_2}$}}}
\put(3676,-1561){\makebox(0,0)[lb]{\smash{\SetFigFont{12}{14.4}{rm}$x$}}}
\put(9076,-1561){\makebox(0,0)[lb]{\smash{\SetFigFont{12}{14.4}{rm}$y$}}}
\put(7276,-1711){\makebox(0,0)[lb]{\smash{\SetFigFont{12}{14.4}{rm}$=$}}}
\put(6376,-1711){\makebox(0,0)[lb]{\smash{\SetFigFont{12}{14.4}{rm}$x\frac{1}{\delta_1}$}}}
\end{picture}

\]

\bs
{\sc Proof.}  Just apply the appropriate annular tangles
and use Reidemeister moves. \qed

\bs Thus we could rewrite equations for $x\in P^u$ entirely in
terms of $x$. The least obvious reformulation of these equations
involves less boundary points than above and requires positivity.

\bs\ni
{\bf Theorem 2.11.8} Let $P$ be a spherical finite-dimensional
$C^*$-planar algebra and $u\in P_2$ be bi-unitary. Then $x\in P^u$
iff
\[
	\begin{picture}(0,0)%
\epsfig{file=ps/pic55.ps}%
\end{picture}%
\setlength{\unitlength}{0.00041700in}%
%
\begingroup\makeatletter\ifx\SetFigFont\undefined
% extract first six characters in \fmtname
\def\x#1#2#3#4#5#6#7\relax{\def\x{#1#2#3#4#5#6}}%
\expandafter\x\fmtname xxxxxx\relax \def\y{splain}%
\ifx\x\y   % LaTeX or SliTeX?
\gdef\SetFigFont#1#2#3{%
  \ifnum #1<17\tiny\else \ifnum #1<20\small\else
  \ifnum #1<24\normalsize\else \ifnum #1<29\large\else
  \ifnum #1<34\Large\else \ifnum #1<41\LARGE\else
     \huge\fi\fi\fi\fi\fi\fi
  \csname #3\endcsname}%
\else
\gdef\SetFigFont#1#2#3{\begingroup
  \count@#1\relax \ifnum 25<\count@\count@25\fi
  \def\x{\endgroup\@setsize\SetFigFont{#2pt}}%
  \expandafter\x
    \csname \romannumeral\the\count@ pt\expandafter\endcsname
    \csname @\romannumeral\the\count@ pt\endcsname
  \csname #3\endcsname}%
\fi
\fi\endgroup
\begin{picture}(8220,3047)(2379,-3084)
\put(3676,-1561){\makebox(0,0)[lb]{\smash{\SetFigFont{12}{14.4}{rm}$x$}}}
\put(9077,-1562){\makebox(0,0)[lb]{\smash{\SetFigFont{12}{14.4}{rm}$x$}}}
\put(5986,-1711){\makebox(0,0)[lb]{\smash{\SetFigFont{12}{14.4}{rm}$=$}}}
\put(7051,-1681){\makebox(0,0)[lb]{\smash{\SetFigFont{12}{14.4}{rm}$\delta^2$}}}
\end{picture}
.
\]

\bs
{\sc Proof.}  $(\Rightarrow)$ This is easy and requires no
positivity.

$(\Leftarrow)$ We begin by observing that
$\d\s^*_u(z)$ is given by the tangle below
\[
	\begin{picture}(0,0)%
\epsfig{file=ps/pic56.ps}%
\end{picture}%
\setlength{\unitlength}{0.00041700in}%
%
\begingroup\makeatletter\ifx\SetFigFont\undefined
% extract first six characters in \fmtname
\def\x#1#2#3#4#5#6#7\relax{\def\x{#1#2#3#4#5#6}}%
\expandafter\x\fmtname xxxxxx\relax \def\y{splain}%
\ifx\x\y   % LaTeX or SliTeX?
\gdef\SetFigFont#1#2#3{%
  \ifnum #1<17\tiny\else \ifnum #1<20\small\else
  \ifnum #1<24\normalsize\else \ifnum #1<29\large\else
  \ifnum #1<34\Large\else \ifnum #1<41\LARGE\else
     \huge\fi\fi\fi\fi\fi\fi
  \csname #3\endcsname}%
\else
\gdef\SetFigFont#1#2#3{\begingroup
  \count@#1\relax \ifnum 25<\count@\count@25\fi
  \def\x{\endgroup\@setsize\SetFigFont{#2pt}}%
  \expandafter\x
    \csname \romannumeral\the\count@ pt\expandafter\endcsname
    \csname @\romannumeral\the\count@ pt\endcsname
  \csname #3\endcsname}%
\fi
\fi\endgroup
\begin{picture}(3644,3046)(2379,-3083)
\put(4426,-1636){\makebox(0,0)[lb]{\smash{\SetFigFont{12}{14.4}{rm}$z$}}}
\end{picture}

\]
\ni Here the adjoint of $\s_u$ is as a map between the
finite-dimensional Hilbert spaces $P_k$ and $P_{k+1}$ with inner
products given by the normalized traces.  This formula for $\s^*_u$
thus follows from the equality by isotopy of the following two
networks, the first of which is, up to a power of $\d$,  \
$\<\s_u(x),z^*\>$
\[
	\begin{picture}(0,0)%
\epsfig{file=ps/pic57.ps}%
\end{picture}%
\setlength{\unitlength}{0.00041700in}%
%
\begingroup\makeatletter\ifx\SetFigFont\undefined
% extract first six characters in \fmtname
\def\x#1#2#3#4#5#6#7\relax{\def\x{#1#2#3#4#5#6}}%
\expandafter\x\fmtname xxxxxx\relax \def\y{splain}%
\ifx\x\y   % LaTeX or SliTeX?
\gdef\SetFigFont#1#2#3{%
  \ifnum #1<17\tiny\else \ifnum #1<20\small\else
  \ifnum #1<24\normalsize\else \ifnum #1<29\large\else
  \ifnum #1<34\Large\else \ifnum #1<41\LARGE\else
     \huge\fi\fi\fi\fi\fi\fi
  \csname #3\endcsname}%
\else
\gdef\SetFigFont#1#2#3{\begingroup
  \count@#1\relax \ifnum 25<\count@\count@25\fi
  \def\x{\endgroup\@setsize\SetFigFont{#2pt}}%
  \expandafter\x
    \csname \romannumeral\the\count@ pt\expandafter\endcsname
    \csname @\romannumeral\the\count@ pt\endcsname
  \csname #3\endcsname}%
\fi
\fi\endgroup
\begin{picture}(9194,6211)(1764,-5783)
\put(2461,-1366){\makebox(0,0)[lb]{\smash{\SetFigFont{12}{14.4}{rm}$z$}}}
\put(2266,-3451){\makebox(0,0)[lb]{\smash{\SetFigFont{12}{14.4}{rm}$x$}}}
\put(8221,-1351){\makebox(0,0)[lb]{\smash{\SetFigFont{12}{14.4}{rm}$z$}}}
\put(8041,-3496){\makebox(0,0)[lb]{\smash{\SetFigFont{12}{14.4}{rm}$x$}}}
\put(6226,-2776){\makebox(0,0)[lb]{\smash{\SetFigFont{12}{14.4}{rm}$=$}}}
\end{picture}

\]
\ni Thus orthogonal projection $E$ onto $\s_u(P_k)$ is
$\s_u\s^*_u$ which is given on $z\in P_k$ by $\frac 1\d$ times the
following picture
\[
	\begin{picture}(0,0)%
\epsfig{file=ps/pic58.ps}%
\end{picture}%
\setlength{\unitlength}{0.00041700in}%
%
\begingroup\makeatletter\ifx\SetFigFont\undefined
% extract first six characters in \fmtname
\def\x#1#2#3#4#5#6#7\relax{\def\x{#1#2#3#4#5#6}}%
\expandafter\x\fmtname xxxxxx\relax \def\y{splain}%
\ifx\x\y   % LaTeX or SliTeX?
\gdef\SetFigFont#1#2#3{%
  \ifnum #1<17\tiny\else \ifnum #1<20\small\else
  \ifnum #1<24\normalsize\else \ifnum #1<29\large\else
  \ifnum #1<34\Large\else \ifnum #1<41\LARGE\else
     \huge\fi\fi\fi\fi\fi\fi
  \csname #3\endcsname}%
\else
\gdef\SetFigFont#1#2#3{\begingroup
  \count@#1\relax \ifnum 25<\count@\count@25\fi
  \def\x{\endgroup\@setsize\SetFigFont{#2pt}}%
  \expandafter\x
    \csname \romannumeral\the\count@ pt\expandafter\endcsname
    \csname @\romannumeral\the\count@ pt\endcsname
  \csname #3\endcsname}%
\fi
\fi\endgroup
\begin{picture}(3644,3046)(2379,-3083)
\put(3676,-1561){\makebox(0,0)[lb]{\smash{\SetFigFont{12}{14.4}{rm}$z$}}}
\end{picture}

\]
\ni Orthogonal projection $F$ onto $P_{1,k+1}$ is given by
$\d F(z)= \quad \boxh{z}$.
An element $w$ of $\s_u(P_k)$ is in $P_{1,k+1}$  iff
$EF(w)=w$. But if $x$ satisfies the condition of the theorem we have
\bs
\[
	\begin{picture}(0,0)%
\epsfig{file=ps/pic59.ps}%
\end{picture}%
\setlength{\unitlength}{0.00041700in}%
%
\begingroup\makeatletter\ifx\SetFigFont\undefined
% extract first six characters in \fmtname
\def\x#1#2#3#4#5#6#7\relax{\def\x{#1#2#3#4#5#6}}%
\expandafter\x\fmtname xxxxxx\relax \def\y{splain}%
\ifx\x\y   % LaTeX or SliTeX?
\gdef\SetFigFont#1#2#3{%
  \ifnum #1<17\tiny\else \ifnum #1<20\small\else
  \ifnum #1<24\normalsize\else \ifnum #1<29\large\else
  \ifnum #1<34\Large\else \ifnum #1<41\LARGE\else
     \huge\fi\fi\fi\fi\fi\fi
  \csname #3\endcsname}%
\else
\gdef\SetFigFont#1#2#3{\begingroup
  \count@#1\relax \ifnum 25<\count@\count@25\fi
  \def\x{\endgroup\@setsize\SetFigFont{#2pt}}%
  \expandafter\x
    \csname \romannumeral\the\count@ pt\expandafter\endcsname
    \csname @\romannumeral\the\count@ pt\endcsname
  \csname #3\endcsname}%
\fi
\fi\endgroup
\begin{picture}(7875,1654)(151,-2800)
\put(151,-2086){\makebox(0,0)[lb]{\smash{\SetFigFont{12}{14.4}{rm}$EF(\sigma_u(x))=\frac{1}{\delta}($}}}
\put(3826,-2011){\makebox(0,0)[lb]{\smash{\SetFigFont{12}{14.4}{rm}$x$}}}
\put(4951,-2086){\makebox(0,0)[lb]{\smash{\SetFigFont{12}{14.4}{rm}$) = \frac{1}{\delta^2}$}}}
\put(6826,-2011){\makebox(0,0)[lb]{\smash{\SetFigFont{12}{14.4}{rm}$x$}}}
\put(8026,-2086){\makebox(0,0)[lb]{\smash{\SetFigFont{12}{14.4}{rm}$=\sigma_u(x)$.}}}
\end{picture}

\]
Hence $x\in P^u$. \qed

\bs {\bf Remark.} We did not use the full force of the hypotheses.
The result will hold in a finite-dimensional general $C^*$-planar
algebra provided isolated circles can be removed with a
multiplicative factor $\d$, that the tangle formula for orthogonal
projection onto $P_{1,k+1}$ is correct, and dim $P_1=1$.

We will see in the case of spin models that $P^u$ may be planar
although $P$ is not. (But the conditions of the above remark are
satisfied by $P^\s$.)

It is easy to check that a bi-invertible $u\in P_2$ may be altered
by four invertible elements $A,B,C,D$ in $P_1$ as in Figure 2.11.9
%%%%%%%%%%%%%%%%%%%%%%%%%%%%%%%%%%%%%%%%%%%%
\begin{figure}
\[
	\begin{picture}(0,0)%
\epsfig{file=ps/pic60.ps}%
\end{picture}%
\setlength{\unitlength}{0.00041700in}%
%
\begingroup\makeatletter\ifx\SetFigFont\undefined
% extract first six characters in \fmtname
\def\x#1#2#3#4#5#6#7\relax{\def\x{#1#2#3#4#5#6}}%
\expandafter\x\fmtname xxxxxx\relax \def\y{splain}%
\ifx\x\y   % LaTeX or SliTeX?
\gdef\SetFigFont#1#2#3{%
  \ifnum #1<17\tiny\else \ifnum #1<20\small\else
  \ifnum #1<24\normalsize\else \ifnum #1<29\large\else
  \ifnum #1<34\Large\else \ifnum #1<41\LARGE\else
     \huge\fi\fi\fi\fi\fi\fi
  \csname #3\endcsname}%
\else
\gdef\SetFigFont#1#2#3{\begingroup
  \count@#1\relax \ifnum 25<\count@\count@25\fi
  \def\x{\endgroup\@setsize\SetFigFont{#2pt}}%
  \expandafter\x
    \csname \romannumeral\the\count@ pt\expandafter\endcsname
    \csname @\romannumeral\the\count@ pt\endcsname
  \csname #3\endcsname}%
\fi
\fi\endgroup
\begin{picture}(1544,4544)(2679,-4583)
\put(4126,-736){\makebox(0,0)[lb]{\smash{\special{ps:gsave currentpoint currentpoint translate
-180.0 rotate neg exch neg exch translate}\SetFigFont{12}{14.4}{rm}$B$\special{ps:currentpoint grestore moveto}}}}
\put(4126,-3436){\makebox(0,0)[lb]{\smash{\special{ps:gsave currentpoint currentpoint translate
-180.0 rotate neg exch neg exch translate}\SetFigFont{12}{14.4}{rm}$C$\special{ps:currentpoint grestore moveto}}}}
\put(3376,-2236){\makebox(0,0)[lb]{\smash{\SetFigFont{12}{14.4}{rm}$u$}}}
\put(2851,-3736){\makebox(0,0)[lb]{\smash{\SetFigFont{12}{14.4}{rm}$D$}}}
\put(2851,-1036){\makebox(0,0)[lb]{\smash{\SetFigFont{12}{14.4}{rm}$A$}}}
\end{picture}

\]
\caption[2.11.9]{}\label{pic60}
\end{figure}
%%%%%%%%%%%%%%%%%%%%%%%%%%%%%%%%%%%%%%%%%
\bs

{\bf Definition 2.11.10.} Two bi-invertibles are said to differ by
a {\it gauge transformation} if one is obtained from the other as
in Figure 2.11.9.

\bs Gauge transformations have an inessential effect on $P^u$; $A$
and $C$ change absolutely nothing, $B$ and $D$ change $P^u$ by a
planar algebra isomorphism (induced by one on $P$-conjugation by
\ $\boxa{d} \ \boxb{d} \ \boxa{d}\dots$).

In the $*$ case, gauge transformations on bi-unitary matrices are
ones with $A,B,C,D$ unitary.

A significant observation about the equations defining $P^u$ above
is that they are {\it linear} so the calculation of $P^u_k$, given
$P_k$ and $u$, is a {\it finite} problem, unlike the calculation of
the planar subalgebra generated by some set, which requires
consideration of infinitely many tangles. In practice, however, the
brute force calculation, even just of dim $P^u_k$, runs into a
serious problem. For the dimension of $P_k$ grows exponentially
with $k$. For $k=2$ the calculation is usually easy enough (indeed
we give an entirely satisfactory general solution for $k=2$ when
$P=P^\s$, below) and somewhat harder for $k=3$. For $k=4$ it tends
to be very demanding even for relatively ``small" $P$'s.
On the other hand, we are dealing with objects  with a lot of
structure. For instance, once we have calculated $P^u_2$ by brute
force or otherwise, the fact that $P^u$ is a planar algebra means
that every unlabelled 2-tangle gives a {\it nonlinear constraint.}
For if a 2-tangle is given labels with elements in  $P$, in order
for the corresponding element of $P_2$ to be in $P^u$, it must lie
in the linear subspace of $P_2$ already calculated. It was the
desire to systematically exploit these highly interesting nonlinear
constraints that led to the theory of planar algebras --- their
generality was only appreciated afterwards.

There are good reasons for wanting to calculate $P^u$. In general
the calculation is greatly facilitated by the presence of group
symmetry but there are many cases of bi-invertible elements with no
apparent symmetry. We hope that the planar algebra $P^u$ plays the
role of ``higher", non-group-like symmetries which reveal
structural properties of the combinatorial object $u$.

We turn now to a special case where this program has had some
partial success, namely in Hadamard matrices.  The theory is no
different for generalized Hadamard matrices, which occur as
biunitaries for spin models. Consider the spin model  $P^\s$ with
$Q$ spins and its spherical planar algebra structure ($\S$2.8). A
bi-invertible element $u$ of $P^\s_2$ is an invertible $Q\x Q$
matrix $u^b_a$ such that
\[
	\begin{picture}(0,0)%
\epsfig{file=ps/pic61.ps}%
\end{picture}%
\setlength{\unitlength}{0.00041700in}%
%
\begingroup\makeatletter\ifx\SetFigFont\undefined
% extract first six characters in \fmtname
\def\x#1#2#3#4#5#6#7\relax{\def\x{#1#2#3#4#5#6}}%
\expandafter\x\fmtname xxxxxx\relax \def\y{splain}%
\ifx\x\y   % LaTeX or SliTeX?
\gdef\SetFigFont#1#2#3{%
  \ifnum #1<17\tiny\else \ifnum #1<20\small\else
  \ifnum #1<24\normalsize\else \ifnum #1<29\large\else
  \ifnum #1<34\Large\else \ifnum #1<41\LARGE\else
     \huge\fi\fi\fi\fi\fi\fi
  \csname #3\endcsname}%
\else
\gdef\SetFigFont#1#2#3{\begingroup
  \count@#1\relax \ifnum 25<\count@\count@25\fi
  \def\x{\endgroup\@setsize\SetFigFont{#2pt}}%
  \expandafter\x
    \csname \romannumeral\the\count@ pt\expandafter\endcsname
    \csname @\romannumeral\the\count@ pt\endcsname
  \csname #3\endcsname}%
\fi
\fi\endgroup
\begin{picture}(9038,1545)(2679,-1734)
\put(4876,-1036){\makebox(0,0)[lb]{\smash{\SetFigFont{12}{14.4}{rm}$=$}}}
\put(9977,-1037){\makebox(0,0)[lb]{\smash{\SetFigFont{12}{14.4}{rm}$=$}}}
\put(2776,-886){\makebox(0,0)[lb]{\smash{\SetFigFont{12}{14.4}{rm}$*$}}}
\put(4276,-1261){\makebox(0,0)[lb]{\smash{\SetFigFont{12}{14.4}{rm}$*$}}}
\put(3001,-1186){\makebox(0,0)[lb]{\smash{\SetFigFont{12}{14.4}{rm}$u$}}}
\put(8176,-1186){\makebox(0,0)[lb]{\smash{\SetFigFont{12}{14.4}{rm}$*$}}}
\put(9076,-886){\makebox(0,0)[lb]{\smash{\SetFigFont{12}{14.4}{rm}$*$}}}
\put(9226,-1111){\makebox(0,0)[lb]{\smash{\SetFigFont{12}{14.4}{rm}$u$}}}
\put(6976,-1036){\makebox(0,0)[lb]{\smash{\SetFigFont{12}{14.4}{rm}and}}}
\put(4531,-736){\makebox(0,0)[lb]{\smash{\special{ps:gsave currentpoint currentpoint translate
-180.0 rotate neg exch neg exch translate}\SetFigFont{12}{14.4}{rm}$u^{-1}$\special{ps:currentpoint grestore moveto}}}}
\put(8461,-736){\makebox(0,0)[lb]{\smash{\special{ps:gsave currentpoint currentpoint translate
-180.0 rotate neg exch neg exch translate}\SetFigFont{12}{14.4}{rm}$u^{-1}$\special{ps:currentpoint grestore moveto}}}}
\end{picture}

\]
\ni So if $(u^{-1})^b_a=v^b_a$ we have $u_{a,b}v_{b,a}=1/Q$ where
the factor $1/Q$ comes from counting oriented circles after
smoothing. If $u$ is biunitary, $v_{a,b}=\overline{u_{b,a}}$, so
the condition is precisely
\[
 |u_{a,b}|=\frac{1}{\sqrt{Q}} %\tag 2.11.11
\]
We call a unitary matrix satisfying 2.11.11 a {\it
generalized Hadamard matrix.} A Hadamard matrix is just $\sqrt{Q}$
times a real generalized Hadamard matrix.

Gauge transformations alter a generalized Hadamard matrix by
multiplying rows and columns by scalars of modulus one ($\pm 1$ in
the Hadamard case). This, together with permutations of the rows
and columns, gives what is called {\it Hadamard equivalence} of
(generalized) Hadamard matrices. Row and column permutations
are easily seen by
2.11.6 to produce equivalent $P^u$'s so any information about $u$
obtained from $P^u$ alone will be {\it invariant under  Hadamard
equivalence.} (The endomorphism $\s_u$ of 2.11.5 itself is more
information than just $P^u$.)

\bs\ni
{\bf Proposition 2.11.12} If $u$ is a generalized Hadamard
matrix, $P^u$ is planar, hence a spherical $C^*$-planar algebra.
Moreover, ${\op{dim}} \ P^u_1=1$, and $P^u_2$ and $P^u_{1,3}$ are
abelian.

\bs
{\sc Proof.} Obviously dim $P^u_1=1$ implies planarity, so
consider a tangle $T$ representing an element of $P_1$. It consists
of a vertical straight line with a 1-box on it, and networks to the
left and right. The networks to the left have exterior shaded white
so only contribute scalars. The picture below is the condition for
such an element to be in  $P^u_1$ (for some element $S$)
\[
	\begin{picture}(0,0)%
\epsfig{file=ps/pic62.ps}%
\end{picture}%
\setlength{\unitlength}{0.00041700in}%
%
\begingroup\makeatletter\ifx\SetFigFont\undefined
% extract first six characters in \fmtname
\def\x#1#2#3#4#5#6#7\relax{\def\x{#1#2#3#4#5#6}}%
\expandafter\x\fmtname xxxxxx\relax \def\y{splain}%
\ifx\x\y   % LaTeX or SliTeX?
\gdef\SetFigFont#1#2#3{%
  \ifnum #1<17\tiny\else \ifnum #1<20\small\else
  \ifnum #1<24\normalsize\else \ifnum #1<29\large\else
  \ifnum #1<34\Large\else \ifnum #1<41\LARGE\else
     \huge\fi\fi\fi\fi\fi\fi
  \csname #3\endcsname}%
\else
\gdef\SetFigFont#1#2#3{\begingroup
  \count@#1\relax \ifnum 25<\count@\count@25\fi
  \def\x{\endgroup\@setsize\SetFigFont{#2pt}}%
  \expandafter\x
    \csname \romannumeral\the\count@ pt\expandafter\endcsname
    \csname @\romannumeral\the\count@ pt\endcsname
  \csname #3\endcsname}%
\fi
\fi\endgroup
\begin{picture}(6806,2744)(2079,-3683)
\put(3001,-2761){\makebox(0,0)[lb]{\smash{\SetFigFont{12}{14.4}{rm}$T$}}}
\put(7858,-1981){\makebox(0,0)[lb]{\smash{\SetFigFont{12}{14.4}{rm}$S$}}}
\put(5158,-2431){\makebox(0,0)[lb]{\smash{\SetFigFont{12}{14.4}{rm}$=$}}}
\end{picture}

\]
\ni If the bottom shaded region is assigned a spin $a$, and the top
region a spin $b$, the left-hand side gives $u^b_aT_a$ and the
right-hand side gives $u^b_aS_b$, so $T_a$ is independent of $a$,
and dim  $P^u_1=1$.  $P^u_{1,3}$ is abelian because $P_{1,3}$ is
and $P^u_2$ is abelian since it is $\s^{-1}_u(\s_u(P_2)\cap P_{1,3})$
by 2.11.6. \qed

\bs So by $\S$4.3, a generalized Hadamard matrix $u$ yields a
subfactor whose planar algebra invariant is $P^u$. In fact such a
subfactor was the starting point of the theory of planar algebras, as the equations for $P^u$ are
those for the relative commutants of a spin model commuting square
given in [JS]. Note that the original subfactor is
hyperfinite whereas the one obtained from 4.3 is not! We now
determine $P^u_2$ for a generalized Hadamard matrix $u$.

\bs {\bf Definition 2.11.13.} Given a $Q\x Q$
generalized Hadamard matrix $u^b_a$ we define the $Q^2\x Q^2$ {\it
profile matrix} \ Prof$(u)$ by
$$
{\op{Prof}}(u)^{c,d}_{a,b} =\sum_x
u^x_a \overline{u^x_b} \ \overline{u^x_c} \ u^x_d \ .
$$
The profile matrix is used in the theory of Hadamard matrices.
We will see that it determines $P^u$.

\bs {\bf Definition 2.11.14.} Given the $Q^2\x Q^2$
matrix \ Prof$(u)$, define the directed graph $\G_u$ on $Q^2$
vertices by $(a,b)\rightarrow\!\!\!- (c,d)$ iff \
Prof$(u)^{c,d}_{a,b}\neq 0$.

\bs The isomorphism class of $\G_u$ is an invariant of Hadamard
equivalence.

\bs\ni
{\bf Theorem 2.11.15} If $u$ is a $Q\x Q$
generalized Hadamard matrix thought of as a biunitary for the spin
model $P^\s$, then the minimal projections of the abelian
$C^*$-algebra $P^u_2$ are in bijection with the connected
components of the graph $\G_u$. Moreover the (normalized) trace of
such a projection is $n/Q^2$ where $n$ is the size of the connected
component, which is necessarily a multiple of $Q$.

\bs
{\sc Proof.} For matrices $x^b_a, \ y^b_a$, the equations of
Theorem 2.11.3 are the ``star-triangle" equations
\[
\sum_d u^d_a \ \overline{u^d_b} \ x^c_d =
u^c_a  \ \overline{u^c_b} \ y^b_a %\tag 2.11.15
\]
which amount to saying that, for each $(a,b)$, the vector
$v_{(a,b)}$ whose $d^{\text{th}}$ component is the
$u^d_a\overline{u^d_b}$ is an eigenvector of the matrix $x^c_d$
with eigenvalue $y^b_a$. The profile matrix is just the matrix of
inner products
$\<v_{(a,b)},v_{(c,d)}\>$ so the orthogonal projection onto the
linear span of $v_{(a,b)}$'s in a connected component is in $P^u_2$
and is necessarily minimal since eigenvectors for distinct minimal
projections are orthogonal.

If the matrix $x$ is an orthogonal projection, $y^b_a$ is either 1
or 0 depending on whether $v_{(a,b)}$ is in the connected
component or not. Consider the picture
\[
	\begin{picture}(0,0)%
\epsfig{file=ps/pic63.ps}%
\end{picture}%
\setlength{\unitlength}{0.00041700in}%
%
\begingroup\makeatletter\ifx\SetFigFont\undefined
% extract first six characters in \fmtname
\def\x#1#2#3#4#5#6#7\relax{\def\x{#1#2#3#4#5#6}}%
\expandafter\x\fmtname xxxxxx\relax \def\y{splain}%
\ifx\x\y   % LaTeX or SliTeX?
\gdef\SetFigFont#1#2#3{%
  \ifnum #1<17\tiny\else \ifnum #1<20\small\else
  \ifnum #1<24\normalsize\else \ifnum #1<29\large\else
  \ifnum #1<34\Large\else \ifnum #1<41\LARGE\else
     \huge\fi\fi\fi\fi\fi\fi
  \csname #3\endcsname}%
\else
\gdef\SetFigFont#1#2#3{\begingroup
  \count@#1\relax \ifnum 25<\count@\count@25\fi
  \def\x{\endgroup\@setsize\SetFigFont{#2pt}}%
  \expandafter\x
    \csname \romannumeral\the\count@ pt\expandafter\endcsname
    \csname @\romannumeral\the\count@ pt\endcsname
  \csname #3\endcsname}%
\fi
\fi\endgroup
\begin{picture}(6432,3604)(1830,-3782)
\put(2218,-1936){\makebox(0,0)[lb]{\smash{\SetFigFont{12}{14.4}{rm}$x$}}}
\put(6343,-1906){\makebox(0,0)[lb]{\smash{\SetFigFont{12}{14.4}{rm}$y$}}}
\put(4738,-2041){\makebox(0,0)[lb]{\smash{\SetFigFont{12}{14.4}{rm}$=$}}}
\end{picture}

\]
\ni
Applying Reidemeister type II moves and summing we obtain the
assertion about the trace. (It is a multiple of $1/Q$ since $x$ is
a $Q\x Q$ matrix.) \qed

\bs If $G$ is a finite abelian group and $g\mapsto\hat g$ is an
isomorphism of $G$ with its dual $\hat G$ (=Hom$(G,\Bbb C^*)$), we
obtain a generalized Hadamard matrix $u$, with $Q=|G|$, by setting
$u^h_g=\frac{1}{\sqrt{Q}} \ \hat h(g)$. We call this a {\it
standard} generalized Hadamard matrix.  It is Hadamard if
$G=(\z/2\z)^n$ for some $n$. We leave it to the reader to check
that if $u$ is standard $P^u$ is exactly the planar
algebra of
$\S$2.9 for the group $G$. In particular, dim$(P^u_k)=Q^k$. It is
well known in subfactor theory that any subfactor with
$N'\cap M_1=\Bbb C^{[M:N]}$ comes from a  group. It can
also be seen directly from association schemes that if
dim$(P^u_2)=Q$ then $u$ is standard up to gauge equivalence
(recall that $P^u$ is always an assocation scheme as remarked in
$\S$2.8).

We have, together with R.~Bacher,
 P.~de la Harpe, and M.G.V.~Bogle performed many computer
calculations. So far we have not found a
generalized Hadamard matrix $u$ for which dim$(P^u_2)=2$ but
dim$(P^u_3) >5$.  Such an example would be a confirmation of our
non-group symmetry program as group-like symmetries tend to show up
in $P_2$. In particular the five $16\x 16$ Hadamard matrices have
dim$ P^u_2=$ 16,8,5,3 and 3, and are completely distinguished by
the trace. There are group-like symmetries in all cases
corresponding to the presence of normalizer in the subfactor
picture.

Haagerup has shown how to construct many interesting examples and
given a complete classification for $Q=5$. In the circulant case he
has shown there are only finitely many examples for fixed prime
$Q$  (see [\qquad ]).

Perhaps somewhat surprisingly, the  presence of a lot of symmetry
in $u$ can cause $P^u_2$ to be small! The kind of biunitary
described in the following result is quite common --- the Paley
type Hadamard matrices give an example.

\bs\ni
{\bf Proposition 2.11.17} Suppose $Q-1$ is prime and let $u$
be a $Q\x Q$ generalized Hadamard matrix with the following two
properties (the first of which is always true up to gauge
equivalence):
\begin{verse}
(i) There is an index $*$ with $u^a_*=u^*_a=1$ for all $a$.

(ii) The group $\z/(Q-1)\z$ acts transitively on the
spins other than $*$, and $u^{ga}_{gb}=u^a_b$ for all
$g\in\z/(Q-1)\z$.
\end{verse}
Then dim$(P^u_2)=2$ or $u$ is gauge equivalent to a standard
matrix.

\bs
{\sc Proof.} The nature of the equations 2.11.15 makes it clear
that $\z/(Q-1)\z$ acts by automorphisms on $P^u_2$, obviously
fixing the projection $e_1$ which is the matrix
$x^b_a=1/Q$. Thus the action preserves $(1-e_1)P^u_2(1-e_1)$.
Since $(Q-1)$ is prime there are  only two possibilities:
either the action is non-trivial and dim$(P^u_2)=Q$ so $P^u$ is
standard, or every solution of 2.11.15 is fixed by
$\z/(Q-1)\z$. In the latter case let $x^b_a,y^b_a$ be a solution
of 2.11.15. Then putting $c=*$ we obtain $\sum_d u^d_a \
\overline{u^d_b} x^*_d=y^b_a$, so $y^b_a$ is determined by the two
numbers
$x^*_*$ and $x^*_d$, $d\neq *$. So by 2.11.7 we are done. \qed

\bs Note that the standard case in the above result can occur. The
$8\x 8$ Hadamard matrix is of the required form, but it is Hadamard
equivalent to a standard matrix. For $Q=$ 12 and 24 this cannot be the
case and dim $P_2^u=2$.

We have very few general results on $P^u_k$ for $k >2$.
We only record the observation that $P^u_k$ is the $\d^2$
eigenspace for the $Q^k\x Q^k$ matrix  given by the ``transfer
matrix with periodic horizontal boundary conditions" for the
$Q$-spin vertex model having the profile matrix as Boltzmann
weights. The transfer matrix is given by the picture:
%%%%%%%%%%%%%%%%%%%%%%%%%%%%%%%%%%%%%%%%%%%%%%%%%%%%%%%%%%
\[
	\begin{picture}(0,0)%
\epsfig{file=ps/pic65.ps}%
\end{picture}%
\setlength{\unitlength}{0.00041700in}%
%
\begingroup\makeatletter\ifx\SetFigFont\undefined
% extract first six characters in \fmtname
\def\x#1#2#3#4#5#6#7\relax{\def\x{#1#2#3#4#5#6}}%
\expandafter\x\fmtname xxxxxx\relax \def\y{splain}%
\ifx\x\y   % LaTeX or SliTeX?
\gdef\SetFigFont#1#2#3{%
  \ifnum #1<17\tiny\else \ifnum #1<20\small\else
  \ifnum #1<24\normalsize\else \ifnum #1<29\large\else
  \ifnum #1<34\Large\else \ifnum #1<41\LARGE\else
     \huge\fi\fi\fi\fi\fi\fi
  \csname #3\endcsname}%
\else
\gdef\SetFigFont#1#2#3{\begingroup
  \count@#1\relax \ifnum 25<\count@\count@25\fi
  \def\x{\endgroup\@setsize\SetFigFont{#2pt}}%
  \expandafter\x
    \csname \romannumeral\the\count@ pt\expandafter\endcsname
    \csname @\romannumeral\the\count@ pt\endcsname
  \csname #3\endcsname}%
\fi
\fi\endgroup
\begin{picture}(8132,3043)(585,-2783)
\end{picture}

\]

\ni where of course the internal spins have been summed over.
This is an immediate consequence of 2.11.8.

We would like to make the following two open problems about
matrices quite explicit. Both concern a generalized Hadamard matrix
$u$.
\begin{verse}
(i) Is the calculation of dim $P^u_k$ feasible in the
polynomial time as a function of $k$?

(ii) Is there a $u$ for which dim
$P^u_k=\frac{1}{k+1}{2k\choose k}$? (i.e.,
$P^u_k$ is just the Temperley-Lieb algebra).
\end{verse}

Finally we make some comments on vertex models. There are many
formal connections with Hopf algebras here which is not surprising
since quantum groups arose from vertex models in statistical
mechanical models ([Dr]). Banica has done some interesting
work from this point of view --- see [Ban].

A vertex model, in the above context, is simply a biunitary (or
biinvertible) in the planar algebra $P^{\otimes}$. The equations of
Theorem 2.11.3 are then just the equations for the higher relative
commutants of a subfactor coming from a certain commuting square
(see [JS]). Perhaps the most interesting examples not coming from
the quantum group machinery are the Krishan-Sunder ``bipermutation
matrices" where $u$ is a permutation matrix with respect to some
basis of the underlying vector space (see [KS]).
B.Bhattacharya has exhibited a planar algebra which is bigger than
that of example 2.3. (Fuss Catalan) and which is necessarily a
planar subalgebra of $P^u$  if $u$ is a bipermutation matrix.




